% ---------------------------------------------------------------------------
% Theorem environments
% ---------------------------------------------------------------------------
\theoremstyle{plain}
\newtheorem{theorem}{Theorem}[section]
\newtheorem{lemma}[theorem]{Lemma}
\newtheorem{corollary}[theorem]{Corollary}
\newtheorem{proposition}[theorem]{Proposition}
\newtheorem{conjecture}[theorem]{Conjecture}
\theoremstyle{definition}
\newtheorem{definition}[theorem]{Definition}
\newtheorem{remark}[theorem]{Remark}
\newtheorem{example}[theorem]{Example}

% The Lean proof is the proof; the text gives a faithful outline of it.
\newenvironment{outline}[1][Outline of the formal proof]{\begin{proof}[#1]}{\end{proof}}
% Results that are not formalized get complete proofs, labelled as such.
\newenvironment{informalproof}{\begin{proof}[Proof (not formalized)]}{\end{proof}}

% ---------------------------------------------------------------------------
% Lean references
%   \lean{name}           a declaration, relative to the namespace Transpose.Tape
%                         unless it starts with "Transpose." (fully qualified)
%   \formal{file}{names}  the formal counterpart of a numbered result: a file
%                         relative to lean/Transpose/ and a comma-separated list
%                         of declarations (same naming convention as \lean)
%   \notformalized        marks a numbered result that has no formal counterpart
% ---------------------------------------------------------------------------
% Lean names are typeset like paths, with line breaks allowed after "_", "." and "/".
\DeclareUrlCommand\leanname{\def\UrlFont{\ttfamily\upshape}}
\renewcommand{\UrlBreaks}{\do\_\do\.\do\/\do\-}
\newcommand{\lean}[1]{\leanname{#1}}
\newcommand{\leanfile}[1]{\leanname{#1}}
\newcommand{\leanlist}[1]{\begingroup\def\leansep{}%
  \renewcommand*{\do}[1]{\leansep\leanname{##1}\def\leansep{, }}\docsvlist{#1}\endgroup}
\newcommand{\formal}[2]{\par\noindent{\small\upshape\raggedright\textsf{Lean:} \leanlist{#2}\unskip
  \nobreak\hfil\penalty50\hskip1em\hbox{}\nobreak\hfill[\leanfile{#1}]\parfillskip=0pt\par}%
  \smallskip\noindent\ignorespaces}
\newcommand{\notformalized}{\par\noindent{\small\upshape\textsf{Not formalized;} a complete proof is given
  below.}\par\smallskip\noindent\ignorespaces}

% ---------------------------------------------------------------------------
% Notation
% ---------------------------------------------------------------------------
\newcommand{\N}{\mathbb{N}}
\newcommand{\Z}{\mathbb{Z}}
\newcommand{\Q}{\mathbb{Q}}
\newcommand{\R}{\mathbb{R}}
\newcommand{\F}{\mathbb{F}}
\newcommand{\T}{\mathsf{T}}
\newcommand{\clog}[1]{\lceil \log_2 #1 \rceil}
\newcommand{\abs}[1]{\lvert #1 \rvert}
\newcommand{\card}[1]{\lvert #1 \rvert}
\newcommand{\Haddr}{\mathcal{H}}
\newcommand{\eone}{e_1}
\newcommand{\etwo}{e_2}
\newcommand{\Mb}{M_{\mathrm{b}}}
\newcommand{\chicross}{\chi_{\mathrm{cross}}}
\newcommand{\src}{\mathrm{src}}
\newcommand{\snk}{\mathrm{snk}}
\newcommand{\eps}{\varepsilon}
\newcommand{\TCC}{\textsc{TCC}}
\newcommand{\dotminus}{\mathbin{\dot{-}}}% truncated subtraction, x \dotminus y = max(x-y, 0)
\DeclareMathOperator{\rank}{rank}
\DeclareMathOperator{\dist}{dist}
\DeclareMathOperator{\poly}{poly}
